\section{Compatibility in the window length}

We next compare the folds at different window lengths.  The results of this
section show that admissible restriction and finite-window normalization are
compatible with one another.

\subsection{Restriction maps}

\begin{definition}[Prefix restrictions]
\label{def:prefix-restriction}
For integers $n\ge m\ge 1$ define
\[
\pi_{n\to m}:X_n\to X_m,
\qquad
\pi_{n\to m}(x_1\cdots x_n):=x_1\cdots x_m.
\]
\end{definition}

\begin{lemma}[Functoriality of admissible restrictions]
\label{lem:pi-functorial}
If $n\ge \ell\ge m\ge 1$, then
\[
\pi_{n\to m}=\pi_{\ell\to m}\circ \pi_{n\to \ell}.
\]
Moreover, for every $c\in \Xfin$ and every $n\ge m$,
\[
\pi_{n\to m}(\pi_n(c))=\pi_m(c).
\]
\end{lemma}

\begin{proof}
All maps involved are literal prefix truncations, so both identities are
immediate.
\end{proof}

\begin{theorem}[Inverse-limit identification]
\label{thm:inverse-limit-golden}
Let
\[
\Xinf
:=
\{c=(c_1,c_2,\ldots)\in\{0,1\}^{\NN}:c_kc_{k+1}=0\text{ for all }k\ge 1\}.
\]
Then there is a natural bijection
\[
\Xinf\cong \varprojlim (X_m,\pi_{m+1\to m}).
\]
\end{theorem}

\begin{proof}
Given $c\in \Xinf$, its prefix family
\[
(\pi_m(c))_{m\ge 1}
\]
is compatible under the restriction maps, hence defines an element of the
inverse limit.  Conversely, given a compatible family
\[
(x^{(m)})_{m\ge 1}\in \varprojlim (X_m,\pi_{m+1\to m}),
\]
define $c_k:=x_k^{(m)}$ for any $m\ge k$.  Compatibility ensures that this is
independent of the choice of $m$.  Since every finite prefix $x^{(m)}$ avoids
$11$, the resulting sequence $c$ lies in $\Xinf$.  The two constructions are
inverse to one another.
\end{proof}

\begin{remark}
The object in Theorem~\ref{thm:inverse-limit-golden} is one sided because the
resolution tower is indexed by prefix length.  It should not be confused with
the two-sided image shift $Y_m$ introduced later at the sequence level.
\end{remark}

\subsection{Compatibility with raw windows}

If one truncates a raw window directly, folding does not in general commute
with restriction.  Let
\[
r_{m+1\to m}:\Omega_{m+1}\to\Omega_m,
\qquad
r_{m+1\to m}(\omega_1\cdots\omega_{m+1})=\omega_1\cdots\omega_m.
\]
For $m=1$ and $\omega=11$ one has
\[
(\Fold_1\circ r_{2\to 1})(11)=\Fold_1(1)=1,
\]
whereas
\[
(\pi_{2\to 1}\circ \Fold_2)(11)=0,
\]
since $N(11)=F_2+F_3=F_4$ and hence $\Fold_2(11)=00$.

The correct raw-level restriction is therefore the folded one.

\begin{definition}[Fold-compatible raw restriction]
\label{def:fold-compatible-restriction}
For $m\ge 1$ define
\[
\widetilde r_{m+1\to m}:\Omega_{m+1}\to\Omega_m,
\qquad
\widetilde r_{m+1\to m}(\omega):=\pi_{m+1\to m}(\Fold_{m+1}(\omega)).
\]
\end{definition}

\begin{proposition}[Commuting diagram for fold-compatible restriction]
\label{prop:fold-omega-commute}
For every $m\ge 1$ and every $\omega\in\Omega_{m+1}$,
\[
\Fold_m(\widetilde r_{m+1\to m}(\omega))
=
\pi_{m+1\to m}(\Fold_{m+1}(\omega)).
\]
Equivalently, the diagram
\[
\begin{CD}
\Omega_{m+1} @>{\Fold_{m+1}}>> X_{m+1}\\
@V{\widetilde r_{m+1\to m}}VV @VV{\pi_{m+1\to m}}V\\
\Omega_m @>{\Fold_m}>> X_m
\end{CD}
\]
commutes.
\end{proposition}

\begin{proof}
By definition,
\[
\widetilde r_{m+1\to m}(\omega)=\pi_{m+1\to m}(\Fold_{m+1}(\omega))\in X_m.
\]
Applying Proposition~\ref{prop:fold-basic}(ii) to this admissible word gives
\[
\Fold_m(\widetilde r_{m+1\to m}(\omega))
=
\widetilde r_{m+1\to m}(\omega)
=
\pi_{m+1\to m}(\Fold_{m+1}(\omega)).
\qedhere
\]
\end{proof}

\begin{remark}[Two raw-level restriction examples]
\label{rem:restriction-examples}
The discrepancy between direct restriction and folded restriction already
appears at the smallest scale, where $m=1$ and $\omega=11$:
\[
(\Fold_1\circ r_{2\to 1})(11)=1,
\qquad
(\Fold_1\circ \widetilde r_{2\to 1})(11)=0.
\]
At the next nontrivial scale one may take $m=3$ and $\omega=1101$.  Then
\[
r_{4\to 3}(1101)=110,
\qquad
\Fold_3(110)=001,
\]
whereas Remark~\ref{rem:two-fold-examples} gives $\Fold_4(1101)=0000$, so
\[
\widetilde r_{4\to 3}(1101)=000,
\qquad
\Fold_3(\widetilde r_{4\to 3}(1101))=000.
\]
Thus the naive truncation retains a visible carry that disappears once one
first folds at the larger resolution and only then restricts.
\end{remark}

